% Mục 4.4. Danh mục thí nghiệm. Người viết: Hiếu.
Bảng~\ref{tab:expl} liệt kê các thí nghiệm của đồ án, mã dùng để gọi chúng trong phần kết quả, mục đích và mục trình bày kết quả của từng thí nghiệm.
Các thí nghiệm chạy tuần tự: chuẩn bị dữ liệu, huấn luyện, chạy từng pipeline để lấy kết quả phát hiện và nhận dạng, rồi mới chạy thí nghiệm.
Mọi số liệu tính trên 176 ảnh kiểm tra, và các bảng trong báo cáo được sinh trực tiếp từ kết quả mà các thí nghiệm ghi ra.

\begin{table}[H]
\centering
\caption{Danh mục thí nghiệm.}\label{tab:expl}
\small
\begin{tabularx}{\linewidth}{@{}l L L c@{}}
\toprule
\textbf{Mã} & \textbf{Tên thí nghiệm} & \textbf{Mục đích} & \textbf{Kết quả} \\
\midrule
E-B1 & Phát hiện không phân lớp & So YOLOv8s với bộ phát hiện cổ điển & Mục~\ref{sec:e-detect} \\
E-B2 & Phát hiện theo lớp & Độ đúng từng lớp của YOLOv8s & Mục~\ref{sec:e-detect} \\
E-C1 & Nhận dạng chữ & So bốn phương án OCR trên vùng cắt theo khung chuẩn & Mục~\ref{sec:e-recognize} \\
E-D1 & Ghép khoản trên nhãn chuẩn & Kiểm tra tính nhất quán của luật ghép & Mục~\ref{sec:e-endtoend} \\
E-D2 & So sánh đầu-cuối & So năm pipeline, gồm cả VLM, trên toàn đơn & Mục~\ref{sec:e-endtoend} \\
E-X1 & Tách nguồn lỗi & Thay lần lượt phát hiện và nhận dạng bằng nhãn chuẩn & Mục~\ref{sec:e-ablation} \\
E-A1 & Độ bền khi ảnh suy giảm & Tác dụng của tiền xử lý ở ba mức ảnh & Mục~\ref{sec:e-ablation} \\
E-A2 & Chỉnh nghiêng & Sai số góc của bước chỉnh nghiêng & Mục~\ref{sec:e-ablation} \\
E-A3 & Nắn phối cảnh và thời gian & Tỉ lệ tìm được tờ giấy và thời gian tiền xử lý & Mục~\ref{sec:e-ablation} \\
\bottomrule
\end{tabularx}
\end{table}
